\documentclass[10pt,a4paper]{amsart}
\usepackage[margin=19mm]{geometry}
\usepackage{amsmath,amssymb,mathtools}
\usepackage[scr=rsfs,cal=euler]{mathalfa}
\usepackage{xfrac}
\usepackage[unicode,hidelinks,bookmarks=false]{hyperref}
\newcommand{\ie}{i.e.\ }
\newcommand{\cf}{cf.\ }
\newcommand{\resp}{respectively\ }
\newcommand{\Subsection}{\subsection}
\providecommand{\nobreakdash}{\nobreak\hbox{-}}
\setlength{\emergencystretch}{2em}
\multlinegap0pt
\usepackage{amssymb}
\usepackage{xcolor}
\newtheorem{lemma}{Lemma}
\newcommand{\ISig}{\mathsf{I}\Sigma^0}
\newcommand{\NN}[0]{\mathbb{N}}
\newcommand{\card}{\operatorname{card}}
\def\qt#1{``#1''}%
\title{The reverse mathematics of\\ bounded Ramsey's theorem for pairs}
\author{Quentin Le Hou\'erou, Ludovic Patey}
\date{Source: arXiv:2509.03688v2 (CC BY 4.0)}
\begin{document}
\maketitle

\noindent\emph{Source and licence.} The reverse mathematics of\\ bounded Ramsey's theorem for pairs. Version arXiv:2509.03688v2 (CC BY 4.0), distributed by arXiv under the Creative Commons Attribution 4.0 International licence, http://creativecommons.org/licenses/by/4.0/, as recorded in the arXiv licence metadata for this exact version and shown on the abstract page https://arxiv.org/abs/2509.03688v2. Full licence notice: \texttt{licenses/arxiv-2509.03688-CC-BY-4.0.txt}.\par\smallskip

\noindent\textit{Editorial scope.} Counted unit: the article's lemma lemma (source0.tex lines 457--464). The article's complete original proof of it is delivered with it (source0.tex lines 466--484, one \texttt{proof} environment of 19 lines). No mathematics is written, repaired or re-derived: the statement and the proof are the article's own lines, in the article's own notation, and no retained line is substituted or re-flowed.  No further source-local context is needed: every cross-reference of every retained line resolves inside the two delivered spans.  Nothing was omitted from any retained span, so the package carries no omission record.  No retained line contains a citation, so the wrapper carries no bibliography and no bibliography entry.  No retained line contains a figure environment, an image-inclusion command or drawing code, so no image is delivered and the package declares no asset dependency.  Editorial formatting only, in addition: an AMS \texttt{amsart} preamble that restates the article's own declarations and macros used by the retained lines, namely the environments \texttt{lemma} on separate counters, together with the standard packages those lines need.  The retained environments print in this document as Lemma 1 in order of appearance, whereas the article prints the same items with the numbering of its own full document.  The complete native source of the article as distributed in the arXiv e-print for this exact version is bundled unchanged as \texttt{sources/184430/source0.tex}.
\begin{lemma}[$\ISig_1$]\label[lemma]{lem:bounded-single-block}
    For every set~$X$, every $b \in \NN$ and $n \in \NN$, there exists a $\Delta_1(X)$ coloring $f_n^b : [\NN]^2 \to 2$ such that:
    \begin{itemize}
        \item There is no set of size $3$ that is $f_n^b$-homogeneous for the color $1$.
        \item For every $e < n$ such that $W^X_e$ is infinite, there exists some $x,y \in W^X_e$ such that $b \leq x < y$ and $f(x,y) = 1$.
    \end{itemize}
    Furthermore, the construction is uniform in $n$ and $b$.
\end{lemma}

\begin{proof}
The coloring $f_n^b$ will be defined by stages, with $f_n^b(x,s)$ being defined at stage $s$ for every $x < s$. At any point of the construction, any $x \in \NN$ will be in state \qt{0} or \qt{1}, meaning that at any further stage~$s$, $f(x, s)$ will be set accordingly. By default, every $x \in \NN$ has the state \qt{0}. We consider a trash $T_s \subseteq [0,s]$ where the elements leaving state \qt{1} will be put, to ensure that they will never go back to that state. Our construction will ensure that $\card T_s \leq n^2$ at every stage $s$.

For every $e < n$, we say that $e$ is \emph{active} at a stage $s$ if $\card (W^X_e[s] \cap [b,s]) > n^2$. Once an element is active, it stays active for the rest of the construction. \\

\textbf{Construction:} at stage $s$, if there exists some new $e < n$ that became active at this stage, we then put every $x < s$ that is in state \qt{1} in state \qt{0} and put them in the trash $T_s$. Then, for every $e' < n$ that is active, we pick an element in $W^X_{e'}[s] \cap [b,s] \setminus T_s$ an put it in state \qt{1}. If no such $e < n$ was found, then the all the elements stay in the same state and no new element is added to the trash.

Finally, whether or not a new $e$ was found, we let $f(x,s) = c$ for every $c < 2$ and every $x < s$ in state \qt{c}.   \\

\textbf{Claim 1}: $T_s \leq n^2$ for every $s \in \NN$, hence the construction is well-defined. Indeed, there are always at most $n$ elements in state \qt{1} at a given time (one per active $e$). And such elements are put into the trash at most $n$ times during the construction (once every time a new $e <n$ became active). Therefore, we can always pick an element in $W^X_{e'}[s] \cap [b,s] \setminus T_s$ if $W^X_{e'}$ is active at stage $s$.
\smallskip

\textbf{Claim 2}: There is no set of size $3$ that is $f_n^b$-homogeneous for the color $1$. Assume by contradiction that we have such a set $\{x < y < z\}$. By construction, $x$ was in state \qt{1} at stage $y$ for $f(x,y)$ to be equal to $1$. Then, for $f(y,z)$ to be equal to $1$, we need $y$ to be in state \qt{1} at stage $z$, but $y$ can only enter state \qt{1} after stage $y$, which would put $x$ in the trash, and make it so that $f(x,z) = 0$.
\smallskip

\textbf{Claim 3}: For every $e < n$ such that $W^X_e$ is infinite, there exists some $x,y \in W^X_e$ such that $b \leq x < y$ and $f(x,y) = 1$.
If $W^X_e$ is infinite, then by $\ISig_1$, there is some stage $s$ such that $\card(W^X_e[s] \cap [b,s]) > n^2$, hence $e$ is eventually active. Then, for every $y > s$, there will be some $x \in W^X_e$ that will be in state \qt{1} at stage $y$, which gives $f(x,y) = 1$. Since $W^X_e$ is infinite, there is one such $y$ in $W^X_e$.

\end{proof}

\end{document}
